\chapter{Introduction}

\section{Background and Motivation}

The persuasiveness of artificial intelligence in generating human-like arguments remains an
open and timely research question. As large language models become increasingly sophisticated,
understanding their effectiveness in argumentation compared to human discourse is critical for
evaluating their deployment in real-world contexts. Yet a largely unexplored dimension concerns
the reverse direction: the extent to which human argumentation can persuade AI systems to shift
their stances and rhetorical approaches. Key questions include: How persuasive are AI-generated
arguments relative to human arguments? What rhetorical patterns and evidence strategies differentiate
AI from human persuasiveness? Can humans persuade AI systems to shift their positions or adapt
their argumentation strategies? How do AI personality configurations moderate both directions of
persuasion? How do participants judge and respond to arguments when they lack information about source?
In this thesis, confirmatory inference focuses on AI-to-human persuasion effects.

Structured debate provides a robust experimental framework for investigating these questions.
In debate contexts, argumentation is structured and directly observable, allowing for
systematic comparison of persuasiveness mechanisms across human and AI agents.
Furthermore, debates produce rich discourse data (argument quality, rhetorical strategies,
participant responses, and AI argument evolution) that can be analyzed to understand the
comparative strengths and weaknesses of AI versus human argumentation, as well as bidirectional
persuasion dynamics wherein human arguments influence AI stance shifts and AI rhetorical strategies
influence human belief change.
A reflection in between debate rounds tests for human ability to reiterate opponent's arguments
and track their own stance evolution.

This study therefore treats persuasion as a dynamic process rather than a single before-after event.
Per-round measurements are used to capture
micro-level updates in stance strength, confidence, and perceived influence from the opponent,
while reflection and acknowledgement entries capture whether participants actually engaged with,
understood, and conceded opponent claims. These process indicators are paired with pre/post
self-reports at debate level to anchor baseline resistance to persuasion and final self-assessed
outcomes. In combination, the design provides temporal resolution (how change unfolds),
mechanistic resolution (why change may occur), and endpoint resolution (how participants interpret
the overall effect), reducing reliance on any single noisy measure.

Information asymmetry in debate scenarios offers a crucial experimental control. By ensuring
participants remain unaware of whether their opponent is human or AI (and what AI model, if
applicable), I can study persuasiveness under reduced source-bias pressure and participant
preconceptions. This setup enables investigation of how participants form judgments about
opponents and reflect on their assumptions and how they shaped the debate.
Opponent identity is withheld until debate completion.

The combination of structured debate, real-time interaction, information-controlled conditions,
and comprehensive data capture creates a controlled opportunity to empirically analyse AI
persuasiveness and contribute evidence-based insights to AI-human argumentation comparison.

\section{Problem Statement}

This thesis investigates AI-mediated persuasion in debate settings, with confirmatory emphasis on
how AI systems influence human belief and confidence outcomes. To conduct this investigation at scale
and under controlled experimental conditions, I design,
implement, and deploy an intelligent debate platform as our primary data collection instrument.
The platform enables collection of rich experimental data---argument content, within-debate belief
changes, AI stance evolution, participant perceptions, and post-debate assessments---suitable for
detailed empirical analysis of AI-to-human persuasion mechanisms in this thesis.
The platform must satisfy the following core requirements:

\begin{enumerate}
    \item \textbf{Controlled Experimental Framework:} Facilitate structured, turn-based live debates
    between human participants and human or AI opponents with consistent
    formatting, timing, and discourse structure to enable fair comparative analysis.

    \item \textbf{Information-Restricted Data Collection:} Implement strict information hiding
    such that participants remain unaware of opponent type (human or AI) and, if AI, which
    specific model or personality generated arguments. This reduces source-bias pressure and enables
    analysis of persuasiveness with lower contamination from participant assumptions. Critically, the platform
    must record AI argument evolution and stance positions across successive debate rounds,
    enabling exploratory analysis of how human arguments may influence AI response modification and
    stance shifts, moderated by personality configuration.

    \item \textbf{Diverse AI Opponent Modeling:} Deploy multiple AI personalities ranging from firm
    to balanced to open-minded with distinct rhetorical styles and evidence-generation strategies,
    all powered by gpt-4o-mini, to enable comparative analysis across different AI argumentation approaches.

    \item \textbf{Comprehensive Data Capture:} Record all debate interactions (arguments, timing,
    meta-data), between-round human user's reflection and belief changes, AI argument positioning and stance
    indicators, and post-debate survey responses to enable downstream confirmatory analysis of AI-to-human
    persuasion and exploratory analysis of argument patterns, stance evolution, and attribution dynamics.

    \item \textbf{Scalable and Reliable Infrastructure:} Design a system capable of orchestrating
    multiple concurrent debates, managing real-time communication, persisting complete debate
    records, and reliably enforcing information barriers throughout the experimental lifecycle.
\end{enumerate}

\section{Research Objectives}

The primary objectives of this project are:

\begin{itemize}
    \item To design and implement a full-stack experimental platform capable of orchestrating
    controlled human-AI and human-human debates and collecting rich discourse data suitable for empirical
    analysis.

    \item To develop an AI argument generation system with multiple configurable personalities
    that produce contextually appropriate, evidence-grounded responses suitable for comparative
    persuasiveness analysis.

    \item To establish and enforce strict information hiding mechanisms that prevent participants
    from inferring opponent type or AI personality during debate, thereby enabling analysis of
    persuasiveness with reduced source-bias and preconception pressure.

    \item To measure persuasion effects in a controlled debate setting, with primary emphasis on
    how distinct AI personalities and rhetorical strategies influence human belief and confidence outcomes,
    under information-hiding conditions that reduce source bias.

    \item To collect comprehensive debate datasets encompassing argument content, discourse
    dynamics, participant reactions, AI stance indicators across rounds, and post-debate survey responses
    to support empirical analysis of AI-to-human persuasion in this thesis.
\end{itemize}

\section{System Architecture and Key Features}

This section provides a high-level overview of the platform architecture and feature set.
Comprehensive implementation details (frontend/backend modules, database schema,
WebSocket specifications, security configuration), including the exact AI
personality prompts (Section~\ref{sec:appendix_ai_personality_prompts}) and the
topic question list (Section~\ref{sec:appendix_topic_questions}), are documented
in Appendix A.

\subsection{Architecture Overview}

The platform is built using a client-server architecture with the following components:

\begin{itemize}
    \item \textbf{Frontend:} A React-based single-page application (SPA) with role-based user
    interfaces for participants and administrators.

    \item \textbf{Backend:} An Express.js server providing REST APIs and WebSocket communication
    for real-time debate interactions.

    \item \textbf{Data Persistence:} MongoDB as the primary data store for users, debates
    (pre/post-debate survey responses, between-round reflections, human stance evolution, AI stance position
    indicators, and complete chat/argument history for bidirectional persuasion analysis).

    \item \textbf{AI Integration:} OpenAI API integration for generating contextually appropriate
    debate arguments based on configurable AI personalities.

    \item \textbf{Real-time Communication:} Socket.IO for bidirectional communication between
    clients and server, enabling real-time argument submission and debate updates.
\end{itemize}

\subsection{Key Features}

\begin{enumerate}
    \item \textbf{Structured Turn-Based Debates:} Debates proceed in defined rounds with
    participants and opponents alternating argument submissions. Each debate maintains a
    complete argument history and discourse context.

    \item \textbf{AI Personality System:} Multiple AI personalities with distinct debating styles
    and instructions can be configured, grounded in the warmth-competence framework from social
    psychology. Arguments are generated using context-aware prompts that include debate history,
    assigned stance, and personality-specific guidance. Distinct personalities (e.g.,
    firm, balanced, open-minded) are designed to systematically vary rhetorical strategies and
    manipulate perceived warmth and competence dimensions, enabling analysis of how different
    communicative styles (emotional openness, rigidity, adaptability) influence AI-to-human persuasion
    in this thesis (prompt templates in Section~\ref{sec:appendix_ai_personality_prompts}).

    \item \textbf{Information Hiding and Controlled Revelation:} Before debate conclusion,
    participants do not observe:
    \begin{itemize}
        \item Whether their opponent is human or AI
        \item Which specific AI model, personality, or system prompt was deployed
        \item Debate game mode, opponent assignment strategy, or experimental condition
    \end{itemize}
    This enables systematic analysis of participant stance evolution in this report by tracking
    how initial positions shift across rounds and culminate in final stances, while preserving
    AI-side traces for future extension studies.

    \item \textbf{Role-Based Access Control:} Separate user interfaces and capabilities for:
    \begin{itemize}
        \item Participants: Limited to viewing their assigned debates (submitting arguments, completing surveys)
        \item Administrators: Access to user management, debate configuration, matchmaking
        dashboards, and system monitoring
    \end{itemize}

    \item \textbf{Automatic Matchmaking:} Scheduled jobs automatically match participants based
    on availability, preference, and system configuration.

    \item \textbf{Session Management:} Persistent session management with JWT-based authentication
    for both registered and guest users, including automatic cleanup of stale sessions.

    \item \textbf{Rate Limiting and Security:} Middleware for rate limiting API requests, CORS
    configuration for secure cross-origin communication, and password hashing for credential
    storage.

    \item \textbf{Post-Debate Assessment:} Comprehensive survey capturing:
    (1) changes in original stance,
    (2) strength and confidence in participant's stance after the debate,
    (3) participant's assumption regarding opponent identity (human, AI, or unsure),
    (4) participant's confidence and reasoned justification for the assumption,
    (5) whether the assumption influenced their persuadability during the debate,
    (6) which range of rounds the participant suspected their opponent's identity during the debate,
    (7) and the detection cues from the argumentation that led to that suspicion.
    This data enables analysis of how assumption accuracy and belief in opponent
    identity correlate with persuasion outcomes.
\end{enumerate}

\section{Scope and Limitations}

\subsection{Scope}

This thesis focuses on:

\begin{itemize}
    \item Architecture and implementation of a full-stack experimental debate platform engineered
    for controlled persuasion data collection
    \item Design and enforcement of information hiding mechanisms to prevent source bias during
    debates
    \item Integration of LLM-based argument generation with configurable personalities and
    context-aware prompts capable of adaptive argument modification
    \item Comprehensive data capture, storage, and retrieval infrastructure for debate datasets including
    human stance evolution/ability to reiterate opponent's arguments and AI argument/stance positioning across debate rounds
    \item Administrative tools for experiment design, personality condition assignment, participant
    management, and controlled debate orchestration
    \item Real-time communication infrastructure ensuring consistent debate flow and reliable recording
    of persuasion indicators
    \item Post-hoc data export and analysis capabilities, with this thesis executing the
    confirmatory AI-to-human persuasion analysis and robustness checks
    \item Inferential scope constrained to AI-to-human persuasion in this thesis;
    human-to-AI persuasion analyses are deferred to future work
\end{itemize}

\subsection{Limitations}

The following aspects are explicitly out of scope or represent known limitations:

\begin{itemize}
    \item Large-scale replication beyond the current pilot sample (descriptive $n=90$; complete-case model $n=89$)
    \item Full argument-quality NLP scoring beyond the current validation scope
    \item Mobile-specific optimizations or native mobile applications
    \item Large-scale distributed deployment or cloud infrastructure considerations
    (single-server reference architecture)
    \item Integration with external recruitment or crowdsourcing platforms
\end{itemize}

\section{Thesis Organization}

The remainder of this thesis is organized as follows:

\begin{itemize}
    \item \textbf{Chapter 2:} Literature Review --- Background on AI persuasion, measurement
    methodologies, and debate-format experimental designs
    \item \textbf{Chapter 3:} System Design and Architecture --- The debate platform as data
    collection infrastructure: design choices, implementation, and data capture mechanisms
    \item \textbf{Chapter 4:} Experimental Methodology --- Study design, participant protocol,
    conditions, and statistical analysis approach
    \item \textbf{Chapter 5:} Results and Analysis --- Experimental findings on AI-to-human
    persuasion (belief and confidence outcomes), robustness checks, and exploratory mechanism analysis
    \item \textbf{Chapter 6:} Conclusion and Future Work --- Synthesis of findings, implications,
    limitations, and directions for future research
\end{itemize}

\section{Contributions}

The primary contributions of this work are:

\begin{enumerate}
    \item Empirical pilot evidence on AI-to-human persuasion: quantitative estimates showing that AI opponents
    with distinct personalities produced belief-change effects whose confirmatory confidence
    intervals include zero at the current sample size.

    \item A methodologically transparent experimental design and information-controlled debate platform
    designed to help isolate AI persuasive effects from source bias, personalization, and structural
    factors, supporting replicable and defensible empirical analysis.

    \item Fine-grained analysis of persuasion dynamics: participant-level belief and confidence trajectories,
    combined with reflection and post-debate attribution data, characterize a directional dissociation pattern
    under non-identified confirmatory effects and examine exploratory mechanism hypotheses.

    \item Evidence on the role of AI personality configuration in persuasion outcomes:
    comparative analysis across firm, balanced, and open-minded personalities indicates no significant
    differences in belief-change outcomes, justifying pooled AI treatment analysis.

    \item A comprehensive, annotated dataset of debate records (arguments, within-debate belief measures,
    attribution data, post-debate assessments) suitable for ongoing and future empirical analysis.
\end{enumerate}
